% Level 2, Week 2 -- Functions That Give Back (student workbook)
\chapterweek{2}{Functions That Give Back}{return hands an answer back; parameters make one function work for many jobs.}

\section{What you will learn}
\begin{itemize}[leftmargin=1.2em]
  \item How \textbf{parameters} let one function do many different jobs.
  \item How \texttt{return} hands an answer back, so your program can use it.
  \item Why printing an answer is not the same as returning it.
  \item How to write a \textbf{docstring} that explains your function.
\end{itemize}

\section{The big idea: a function is a machine}
A function is like a machine. Values go \textbf{in} through the parameters. The machine works on them.
An answer comes \textbf{back out} with \texttt{return}.

\begin{center}
%% alt: A function takes values in, does its work and returns a result. In: the arguments 12 and 5. Function: damage(attack, shield) works out result = attack - shield. Out: the returned value 7.
\begin{tikzpicture}[font=\small,
  box/.style={draw=black,thick,rounded corners=3pt,align=center,inner sep=6pt}]
  \node[box,fill=kcsky!20] (in) {\textbf{IN} (arguments)\\\texttt{12}\quad \texttt{5}};
  \node[box,fill=kcorange!25,right=1.4cm of in,minimum width=4.4cm] (m) {\textbf{FUNCTION}\\\texttt{damage(attack, shield)}\\\texttt{result = attack - shield}};
  \node[box,fill=kcgreen!20,right=1.9cm of m] (out) {\textbf{OUT} (return)\\\texttt{7}};
  \draw[-{Stealth[length=3mm]},thick] (in) -- (m);
  \draw[-{Stealth[length=3mm]},thick] (m) -- node[above]{\texttt{return}} (out);
\end{tikzpicture}
\end{center}

\begin{pycode}[A function with two parameters that returns an answer]
def damage(attack, shield):
    """Return the damage a hit does."""
    return attack - shield

hit = damage(12, 5)       # 12 goes into attack, 5 goes into shield
print(f"Ouch! {hit} damage")
\end{pycode}
The values in the parentheses when you \emph{call} the function (\texttt{12} and \texttt{5}) are \textbf{arguments}.
They fill the \textbf{parameters} (\texttt{attack} and \texttt{shield}) in the same order.

\subsection{Print or return?}
\begin{center}
\renewcommand{\arraystretch}{1.3}
\begin{tabular}{@{}p{0.45\linewidth}p{0.45\linewidth}@{}}
\toprule
\textbf{\texttt{print(answer)}} & \textbf{\texttt{return answer}} \\ \midrule
Shows the answer on the screen. & Hands the answer back to the code that called the function. \\
Like shouting it out of a window: you cannot keep it. & Like putting it in your hand: you can store it or do math with it. \\
The function gives back \texttt{None} (nothing). & The function gives back the answer. \\ \bottomrule
\end{tabular}
\end{center}

\textbf{Bug hunt clue:} if you see the word \texttt{None} where a number should be, a function probably \emph{printed}
instead of \emph{returning}.

\subsection{Docstrings: notes for future-you}
A \textbf{docstring} is a sentence in triple quotes just under the \texttt{def} line. It says what the function gives back.
Thonny shows it when you hover over the function name, and your friends can read it to understand your code.

\subsection{What happens in a function stays in a function}
Variables you make inside a function (like \texttt{result}) only exist inside it. Outside the function, Python says
\texttt{NameError}. If you need the value outside, \texttt{return} it.
\begin{pycode}[Inside stays inside]
def damage(attack, shield):
    result = attack - shield    # result lives only inside damage()
    return result

hit = damage(12, 5)
print(hit)        # 7     the returned value, caught in hit
print(result)     # NameError: name 'result' is not defined
\end{pycode}

\begin{exercise}[Game Math]
Folder: \texttt{exercises/week02/ex-game-math}. Open \texttt{starter.py} and press \menu{Run}.
\begin{enumerate}[leftmargin=1.4em]
  \item Find the word \texttt{None} in the output. Which function caused it? \rule{4cm}{0.4pt}
  \item TODO 1: change \texttt{print(...)} to \texttt{return ...} in \texttt{xp\_needed()}.
  \item Predict first, then finish \texttt{damage()} (TODO 2 and 3):\quad
        \texttt{damage(12, 5)} = \rule{1.5cm}{0.4pt}\qquad \texttt{damage(3, 5)} = \rule{1.5cm}{0.4pt}
  \item TODO 4: finish \texttt{is\_critical()}. Tip: \texttt{return roll >= 19} gives back True or False by itself.
  \item TODO 5: remove the \texttt{\#} from the two health lines. Health left: \rule{1.5cm}{0.4pt}
  \item Check your work with \texttt{check\_my\_work.py}.
\end{enumerate}
\end{exercise}

\begin{exercise}[Flower Garden]
Folder: \texttt{exercises/week02/ex-flower-garden}. Press \menu{Run}: all three flowers land in the same spot with 4 petals.
\begin{enumerate}[leftmargin=1.4em]
  \item TODO 1: \texttt{petal\_turn(petals)} should return \texttt{360 / petals}. How far does it turn for 6 petals? \rule{1.5cm}{0.4pt}
  \item TODO 2 and 3: use \texttt{t.goto(x, y)} so each flower appears where it was asked to.
  \item TODO 4: repeat once for each petal: \texttt{range(petals)}.
  \item Look at the \texttt{def} line: \texttt{size=60} is a \textbf{default value}. If a call leaves out the size, it is 60;
        \texttt{draw\_flower(t, 0, 80, 8, SKY, 70)} hands in 70 instead.
  \item Fill in the arguments you used for your favorite flower:\par
        \texttt{draw\_flower(t,} x = \rule{1.2cm}{0.4pt}, y = \rule{1.2cm}{0.4pt}, petals = \rule{1.2cm}{0.4pt}, color = \rule{2.2cm}{0.4pt}\texttt{)}
\end{enumerate}
\end{exercise}

\begin{vocab}
  \vocabword{parameter}{A named slot in a function's \texttt{def} line, like \texttt{attack}. It receives a value.}
  \vocabword{argument}{The value you hand in when you call the function, like \texttt{12}.}
  \vocabword{return}{Hands an answer back to the code that called the function, and ends the function.}
  \vocabword{None}{Python's word for \enquote{nothing}. A function with no \texttt{return} gives back \texttt{None}.}
  \vocabword{docstring}{A sentence in triple quotes under \texttt{def} that explains what the function does.}
  \vocabword{scope}{Where a variable can be seen. Variables made inside a function stay inside it.}
\end{vocab}

\begin{reflection}
Explain print versus return to someone who missed today (use your own example):
\answerlines{2}
Write the docstring you are most proud of:
\answerlines{1}
\end{reflection}
